\documentclass[11pt]{article}

\usepackage[utf8]{inputenc}
\usepackage[T1]{fontenc}
\usepackage[spanish,es-nodecimaldot]{babel}
\usepackage{amsmath,amssymb,mathtools}
\usepackage{geometry}
\usepackage{booktabs}
\usepackage{array}
\usepackage{enumitem}
\usepackage{hyperref}
\usepackage{xcolor}
\usepackage{siunitx}
\usepackage{graphicx}
\usepackage{float}

\geometry{margin=2.5cm}
\hypersetup{
    colorlinks=true,
    linkcolor=blue!50!black,
    urlcolor=blue!50!black
}

\title{Plataforma Antivibratoria Activa para Rechazo de Perturbaciones\\Version Resumida con IMC-PID}
\author{Documento base para Overleaf}
\date{21 de abril de 2026}

\begin{document}

\maketitle
\tableofcontents
\newpage

\section{Contexto y objetivo}

En el laboratorio se realizan pruebas sobre papas usando un sensor de ultrasonido. La medicion se ve afectada por vibraciones de la mesa o base donde esta montado el sensor. Si la base vibra, la posicion del sensor cambia y la distancia medida deja de reflejar unicamente la geometria real del objeto.

Por esta razon, el problema se plantea como \textbf{rechazo de perturbaciones}. El objetivo de control es que la plataforma donde esta montado el sensor permanezca lo mas inmovil posible aunque la base vibre.

\section{Entradas, salida y orden del sistema}

\begin{center}
\small
\begin{tabular}{>{\raggedright\arraybackslash}p{3.0cm} >{\raggedright\arraybackslash}p{2.8cm} >{\raggedright\arraybackslash}p{7.2cm}}
\toprule
\textbf{Categoria} & \textbf{Simbolo} & \textbf{Interpretacion} \\
\midrule
Entrada de control & $u(t)$ & Voltaje aplicado al actuador. Es la unica entrada de la planta nominal. \\
Perturbacion externa & $z(t)$ o $\ddot z(t)$ & Vibracion de la base. Actua sobre el sistema fisico, pero no se usa como segunda entrada al derivar la funcion de transferencia nominal. \\
Salida regulada & $y(t)=x(t)$ & Posicion absoluta de la plataforma donde esta montado el sensor. \\
Variable interna & $q(t)=x(t)-z(t)$ & Desplazamiento relativo entre plataforma y base. \\
Estados & $q,\ \dot q,\ i$ & Deformacion relativa, velocidad relativa y corriente del actuador. \\
Orden nominal & $3$ & La planta nominal es de tercer orden porque necesita tres estados independientes. \\
\bottomrule
\end{tabular}
\end{center}

La aclaracion importante es:
\begin{quote}
\textbf{la entrada manipulada es el voltaje $u(t)$, mientras que la corriente $i(t)$ es un estado interno del actuador.}
\end{quote}

\section{Modelo fisico}

\subsection{Ecuacion mecanica}

Se define la coordenada relativa:
\begin{equation}
q(t)=x(t)-z(t).
\end{equation}

Aplicando la segunda ley de Newton a la plataforma:
\begin{equation}
m\ddot x = -k(x-z)-c(\dot x-\dot z)+F_a.
\end{equation}

Pasando a coordenada relativa:
\begin{equation}
m\ddot q + c\dot q + kq = F_a - m\ddot z.
\label{eq:mech_full}
\end{equation}

\subsection{Modelo del actuador}

Se asume un actuador electromecanico lineal:
\begin{equation}
F_a = K_f i.
\label{eq:fa_ki}
\end{equation}

La ecuacion electrica del actuador es:
\begin{equation}
L\dot i + Ri + K_e\dot q = u.
\label{eq:elec_full}
\end{equation}

\subsection{Modelo completo}

Sustituyendo \eqref{eq:fa_ki} en \eqref{eq:mech_full}:
\begin{equation}
m\ddot q + c\dot q + kq = K_f i - m\ddot z,
\label{eq:full1}
\end{equation}
\begin{equation}
L\dot i + Ri + K_e\dot q = u.
\label{eq:full2}
\end{equation}

Estas dos ecuaciones representan la planta fisica completa. La perturbacion aparece en el modelo fisico, pero para obtener la \textbf{planta nominal} de control se fija temporalmente en cero.

\section{Planta nominal en espacio de estados y funcion de transferencia}

\subsection{Planta nominal}

Para obtener la planta nominal se toma:
\[
z(t)=0,\qquad \dot z(t)=0,\qquad \ddot z(t)=0.
\]

Entonces \eqref{eq:full1}--\eqref{eq:full2} quedan:
\begin{equation}
m\ddot q + c\dot q + kq = K_f i,
\end{equation}
\begin{equation}
L\dot i + Ri + K_e\dot q = u.
\end{equation}

Se definen los estados:
\[
x_1=q,\qquad x_2=\dot q,\qquad x_3=i.
\]

Por tanto:
\begin{equation}
\dot x = Ax+Bu,
\end{equation}
\begin{equation}
y = Cx+Du,
\end{equation}
con
\begin{equation}
A=
\begin{bmatrix}
0 & 1 & 0\\
-\dfrac{k}{m} & -\dfrac{c}{m} & \dfrac{K_f}{m}\\
0 & -\dfrac{K_e}{L} & -\dfrac{R}{L}
\end{bmatrix},
\qquad
B=
\begin{bmatrix}
0\\
0\\
\dfrac{1}{L}
\end{bmatrix},
\end{equation}
\begin{equation}
C=
\begin{bmatrix}
1 & 0 & 0
\end{bmatrix},
\qquad
D=0.
\end{equation}

\subsection{Funcion de transferencia nominal}

La funcion de transferencia nominal se obtiene con:
\begin{equation}
G(s)=C(sI-A)^{-1}B+D.
\label{eq:ss2tf}
\end{equation}

Al desarrollar \eqref{eq:ss2tf} se obtiene:
\begin{equation}
G(s)=\frac{Q(s)}{U(s)}
=
\frac{K_f}
{(Ls+R)(ms^2+cs+k)+K_fK_e s}.
\label{eq:tf_3rd}
\end{equation}

Expandiendo:
\begin{equation}
G(s)=
\frac{K_f}
{Lm\,s^3 + (Lc+Rm)s^2 + (Lk+Rc+K_fK_e)s + Rk}.
\end{equation}

Esta es la planta nominal de \textbf{tercer orden}. La perturbacion no se incorpora aqui como segunda entrada; simplemente se anula para construir la TF principal.

\section{Modelo reducido para sintonia IMC-PID}

Para sintonizar un PID de forma simple y defendible, se toma una reduccion de orden del actuador asumiendo dinamica electrica rapida:
\begin{equation}
L\dot i \approx 0.
\end{equation}

Entonces:
\begin{equation}
Ri + K_e\dot q \approx u.
\end{equation}

Si se trabaja con un modelo reducido normalizado para sintonia, se usa:
\begin{equation}
G_r(s)=\frac{1}{ms^2+cs+k}.
\label{eq:reduced_tf}
\end{equation}

Con los datos dados:
\[
m=1,\qquad c=5,\qquad k=100,
\]
queda:
\begin{equation}
G_r(s)=\frac{1}{s^2+5s+100}.
\label{eq:tf_num}
\end{equation}

Este es el modelo que se usa para obtener el PID por IMC. Por tanto:

\begin{itemize}[leftmargin=1.2cm]
\item la \textbf{planta fisica nominal} sigue siendo de tercer orden;
\item el \textbf{modelo de sintonia del PID} es de segundo orden.
\end{itemize}

\section{Diseno del controlador IMC-PID}

\subsection{Objetivo nominal}

Se toma como modelo deseado de lazo cerrado:
\begin{equation}
T_d(s)=\frac{1}{\alpha s+1},
\label{eq:td}
\end{equation}
donde $\alpha$ es el parametro de sintonia IMC.

\subsection{Obtencion del controlador}

Usando sintonia IMC:
\begin{equation}
C(s)=\frac{T_d(s)}{G_r(s)\left[1-T_d(s)\right]}.
\end{equation}

Sustituyendo \eqref{eq:tf_num} y \eqref{eq:td}:
\begin{equation}
C(s)=
\frac{\dfrac{1}{\alpha s+1}}
{\dfrac{1}{s^2+5s+100}\left(1-\dfrac{1}{\alpha s+1}\right)}.
\end{equation}

Como
\[
1-\frac{1}{\alpha s+1}=\frac{\alpha s}{\alpha s+1},
\]
se obtiene:
\begin{equation}
C(s)=\frac{s^2+5s+100}{\alpha s}.
\label{eq:c_alpha}
\end{equation}

Separando terminos:
\begin{equation}
C(s)=\frac{1}{\alpha}s+\frac{5}{\alpha}+\frac{100}{\alpha s}.
\label{eq:c_parallel}
\end{equation}

\subsection{Forma PID ideal}

La forma ideal es:
\begin{equation}
C_{PID}(s)=K_c\left(1+\frac{1}{\tau_I s}+\tau_D s\right).
\label{eq:pid_ideal}
\end{equation}

Expandiendo \eqref{eq:pid_ideal}:
\begin{equation}
C_{PID}(s)=K_c+\frac{K_c}{\tau_I}\frac{1}{s}+K_c\tau_D s.
\label{eq:pid_ideal_exp}
\end{equation}

Ahora se comparan \eqref{eq:c_parallel} y \eqref{eq:pid_ideal_exp} termino a termino:

\begin{equation}
K_c=\frac{5}{\alpha},
\label{eq:kc_alpha}
\end{equation}
\begin{equation}
\frac{K_c}{\tau_I}=\frac{100}{\alpha},
\end{equation}
\begin{equation}
K_c\tau_D=\frac{1}{\alpha}.
\end{equation}

Usando \eqref{eq:kc_alpha}, se obtiene:
\begin{equation}
\tau_I=\frac{K_c}{100/\alpha}
=
\frac{5/\alpha}{100/\alpha}
=
\frac{5}{100}
=
0.05\ \text{s},
\end{equation}
\begin{equation}
\tau_D=\frac{(1/\alpha)}{K_c}
=
\frac{1/\alpha}{5/\alpha}
=
\frac{1}{5}
=
0.2\ \text{s}.
\end{equation}

Por tanto:
\begin{equation}
K_c=\frac{5}{\alpha},\qquad
\tau_I=0.05\ \text{s},\qquad
\tau_D=0.2\ \text{s}.
\label{eq:pid_params_general}
\end{equation}

\subsection{Mantener $P=50$}

Si se desea que la ganancia proporcional del PID ideal sea:
\[
P=K_c=50,
\]
entonces, usando \eqref{eq:kc_alpha}:
\begin{equation}
50=\frac{5}{\alpha}.
\end{equation}

Despejando:
\begin{equation}
\alpha=\frac{5}{50}=0.1\ \text{s}.
\end{equation}

Por tanto, con $P=50$:
\begin{equation}
K_c=50,\qquad
\tau_I=0.05\ \text{s},\qquad
\tau_D=0.2\ \text{s}.
\label{eq:pid_final}
\end{equation}

El controlador PID ideal final queda:
\begin{equation}
\boxed{
C(s)=50\left(1+\frac{1}{0.05\,s}+0.2\,s\right)
}
\label{eq:pid_ideal_final}
\end{equation}

\subsection{Forma paralela equivalente}

Expandiendo \eqref{eq:pid_ideal_final}:
\begin{equation}
C(s)=50+\frac{1000}{s}+10s.
\end{equation}

Por tanto:
\[
K_p=50,\qquad K_i=1000,\qquad K_d=10.
\]

\section{Equivalencia con el bloque PID de Simulink}

El bloque que se quiere usar tiene la forma:
\begin{equation}
C_b(s)=P\left(1+\frac{1}{Is}+D\frac{N}{1+N/s}\right).
\label{eq:block_pid}
\end{equation}

\subsection{Paso 1: comparar la parte proporcional e integral}

Comparando \eqref{eq:pid_ideal_final} con \eqref{eq:block_pid}, las dos primeras partes son directas:

\begin{equation}
P=K_c=50,
\end{equation}
\begin{equation}
I=\tau_I=0.05,
\end{equation}
\begin{equation}
D=\tau_D=0.2.
\end{equation}

\subsection{Paso 2: entender el termino derivativo filtrado}

En el PID ideal, el termino derivativo es:
\[
D\,s.
\]

En el bloque, el termino derivativo se implementa como:
\[
D\frac{N}{1+N/s}.
\]

Multiplicando numerador y denominador por $s$:
\begin{equation}
D\frac{N}{1+N/s}
=
D\frac{Ns}{s+N}.
\label{eq:filtered_derivative}
\end{equation}

Si $N$ es grande, entonces:
\[
\frac{Ns}{s+N}\approx s,
\]
y por tanto:
\[
D\frac{Ns}{s+N}\approx Ds.
\]

Eso significa que el bloque no usa una derivada ideal pura, sino una derivada \textbf{filtrada}, lo cual es deseable porque evita amplificar ruido.

\subsection{Paso 3: valores para el bloque}

La equivalencia exacta de parametros es:
\[
P=50,\qquad I=0.05,\qquad D=0.2.
\]

El parametro $N$ no sale del PID ideal; se escoge para el filtro derivativo. Un valor inicial razonable es:
\[
N=10.
\]

Entonces, en el bloque de Simulink se pueden poner:
\begin{equation}
\boxed{
P=50,\qquad I=0.05,\qquad D=0.2,\qquad N=10
}
\label{eq:block_values}
\end{equation}

\section{Resultados y figuras}

En esta seccion se dejan los espacios para insertar las graficas y describir los resultados mas importantes.

\subsection{Lazo abierto teorico con funcion de transferencia}

Aqui se debe mostrar la respuesta al escalon del modelo teorico reducido:
\[
G_r(s)=\frac{1}{s^2+5s+100}.
\]

\begin{figure}[H]
\centering
\includegraphics[width=0.82\textwidth]{fig_lazo_abierto_tf.png}
\caption{Respuesta en lazo abierto del modelo teorico obtenido a partir de la funcion de transferencia.}
\label{fig:abierto_tf}
\end{figure}

\noindent\textbf{Texto sugerido para completar:} En la Fig.~\ref{fig:abierto_tf} se presenta la respuesta en lazo abierto del modelo teorico. Aqui se deben reportar tiempo de subida, tiempo de establecimiento, sobreimpulso y valor final.

\subsection{Lazo abierto Simscape}

En esta parte se debe mostrar la respuesta del modelo realista en Simscape bajo la misma excitacion usada en el modelo teorico, para comparar la planta aproximada con la planta fisica.

\begin{figure}[H]
\centering
\includegraphics[width=0.82\textwidth]{fig_lazo_abierto_simscape.png}
\caption{Respuesta en lazo abierto del modelo realista implementado en Simscape.}
\label{fig:abierto_simscape}
\end{figure}

\noindent\textbf{Texto sugerido para completar:} En la Fig.~\ref{fig:abierto_simscape} se observa la respuesta del modelo realista en lazo abierto. Aqui se comparan diferencias con respecto al modelo teorico, especialmente en rapidez, amortiguamiento, saturaciones o efectos no ideales.

\subsection{Lazo cerrado Simscape}

En esta parte se debe mostrar la respuesta del sistema realista en lazo cerrado con el controlador IMC-PID. La referencia natural del problema es cero, y la perturbacion se introduce a traves del movimiento de la base.

\begin{figure}[H]
\centering
\includegraphics[width=0.82\textwidth]{fig_lazo_cerrado_simscape.png}
\caption{Respuesta en lazo cerrado del modelo realista en Simscape con el controlador IMC-PID.}
\label{fig:cerrado_simscape}
\end{figure}

\noindent\textbf{Texto sugerido para completar:} En la Fig.~\ref{fig:cerrado_simscape} se muestra el desempeno del controlador IMC-PID frente a perturbaciones de base. Aqui se debe analizar si el controlador reduce la transmision de vibracion hacia la plataforma y si mantiene el sensor suficientemente estable para la medicion.

\subsection{Tabla de metricas}

\begin{table}[H]
\centering
\small
\begin{tabular}{lcccc}
\toprule
\textbf{Caso} & \textbf{$t_r$} & \textbf{$t_s$} & \textbf{$M_p$} & \textbf{$e_{ss}$} \\
\midrule
Lazo abierto teorico & -- & -- & -- & -- \\
Lazo abierto Simscape & -- & -- & -- & -- \\
Lazo cerrado Simscape & -- & -- & -- & -- \\
\bottomrule
\end{tabular}
\caption{Metricas de desempeno a completar con los resultados obtenidos en simulacion.}
\label{tab:metricas}
\end{table}

\section{Conclusiones}

El problema se formul\'o como rechazo de perturbaciones para estabilizar un sensor de ultrasonido montado sobre una plataforma sensible a vibraciones. La planta fisica nominal se model\'o como un sistema de tercer orden, mientras que para la sintonia del controlador se utiliz\'o un modelo reducido de segundo orden. Con los parametros $m=1$, $c=5$ y $k=100$, y manteniendo $P=50$, se obtuvo un controlador IMC-PID ideal equivalente a:
\[
C(s)=50\left(1+\frac{1}{0.05\,s}+0.2\,s\right).
\]

La forma equivalente para el bloque PID con filtro derivativo es:
\[
P=50,\qquad I=0.05,\qquad D=0.2,\qquad N=10.
\]

El siguiente paso del informe consiste en insertar las figuras de simulacion y completar la tabla de metricas con los resultados del modelo teorico y del modelo realista en Simscape.

\end{document}
